\documentclass[10pt,a4paper]{amsart}
\usepackage[margin=19mm]{geometry}
\usepackage{amsmath,amssymb,mathtools}
\usepackage[scr=rsfs,cal=euler]{mathalfa}
\usepackage{xfrac}
\usepackage[unicode,hidelinks,bookmarks=false]{hyperref}
\newcommand{\ie}{i.e.\ }
\newcommand{\cf}{cf.\ }
\newcommand{\resp}{respectively\ }
\newcommand{\Subsection}{\subsection}
\providecommand{\nobreakdash}{\nobreak\hbox{-}}
\setlength{\emergencystretch}{2em}
\multlinegap0pt
\usepackage{xcolor}
\usepackage{amssymb}
\usepackage{enumerate}
\usepackage{tikz}
\newtheorem{proposition}{Proposition}
\def\ov{\overline}
\def\ss{\sigma}
\title{Homology of matching complexes and representations of symmetric groups}
\author{Michael Bate, Brent Everitt, Sam Ford, Eric Ramos}
\date{Source: arXiv:2312.13750v3 (CC BY 4.0)}
\begin{document}
\maketitle

\noindent\emph{Source and licence.} Homology of matching complexes and representations of symmetric groups. Version arXiv:2312.13750v3 (CC BY 4.0), distributed by arXiv under the Creative Commons Attribution 4.0 International licence, http://creativecommons.org/licenses/by/4.0/, as recorded in the arXiv licence metadata for this exact version and shown on the abstract page https://arxiv.org/abs/2312.13750v3. Full licence notice: \texttt{licenses/arxiv-2312.13750-CC-BY-4.0.txt}.\par\smallskip

\noindent\textit{Editorial scope.} Counted unit: the article's proposition prop1 (source0.tex lines 551--554). The article's complete original proof of it is delivered with it (source0.tex lines 556--589, one \texttt{proof} environment of 34 lines). No mathematics is written, repaired or re-derived: the statement and the proof are the article's own lines, in the article's own notation, and no retained line is substituted or re-flowed.  Retained uncounted as the source-local context the proof needs: lines 496--501.  Nothing was omitted from any retained span, so the package carries no omission record.  No retained line contains a citation, so the wrapper carries no bibliography and no bibliography entry.  No retained line contains a figure environment, an image-inclusion command or drawing code, so no image is delivered and the package declares no asset dependency.  Editorial formatting only, in addition: an AMS \texttt{amsart} preamble that restates the article's own declarations and macros used by the retained lines, namely the environments \texttt{proposition} on separate counters, together with the standard packages those lines need.  The retained environments print in this document as Proposition 1 in order of appearance, whereas the article prints the same items with the numbering of its own full document.  The complete native source of the article as distributed in the arXiv e-print for this exact version is bundled unchanged as \texttt{sources/151885/source0.tex}.
\begin{equation}
\label{eq5}
    \overline{\sigma}_t
\cap
{\textstyle \bigcup}_{i=1}^{t-1}\,\overline{\sigma}_i,
\end{equation}

\begin{proposition}
  \label{prop1}
This order is a shelling of $X(n)$.
\end{proposition}

\begin{proof}
Let $\ss_t=x_0\ldots x_k z_1\ldots z_\ell$ be a facet. We claim that
the subcomplex (\ref{eq5}) is the union of all the simplices
$\overline{\tau}$ where:
\begin{equation}
  \label{eq:2}
\tau=x_0\ldots \widehat{x}_i \ldots x_k z_1\ldots z_\ell,
\end{equation}
with $0\leq i\leq k$ and the hat denoting omission. As each such
$\overline{\tau}$ is a $(\dim \sigma_t-1)$-dimensional sub-simplex of
$\overline{\sigma}$, the proposition 
follows. To see the claim, if in (\ref{eq:2}) the block
$x_i=\{x_{i1},\ldots,x_{ij}\}$, then $\tau$ is a face of the facet:
$$
\ss_s=x_0\ldots \widehat{x}_i \ldots x_k z_1\ldots z_\ell x_{i1}\ldots
x_{ij},
$$
and $\ss_s<\ss_t$. Thus $\ov{\tau}\subseteq\ov{\ss}_s\in
  \bigcup_{i=1}^{t-1}\,\overline{\sigma}_i$ and $\ov{\tau}$ is indeed
  contained in the subcomplex (\ref{eq5}).
%
On the other hand, let $\omega$ be a face of $\ss_t$ of the form:
$$
\omega=x_0\ldots x_k \widehat{z}_1\ldots\widehat{z}_jz_{j+1}\ldots z_\ell.
$$
If $\ss_s$ is a facet that contains $\omega$ as a face then
$\ss_s=
x_0\ldots x_k y_1\ldots y_r z_{j+1}\ldots z_\ell$
where $y_1\ldots y_r$ is some partition of the set
$\{z_1,\ldots,z_j\}$. In particular, $\ss_s\geq \ss_t$ and
$\ov{\omega}\not\in \bigcup_{i=1}^{t-1}\,\overline{\sigma}_i$. This
proves the claim. 
  \qed
\end{proof}

\end{document}
